
Далее мы приведём похожий на алгоритм Дейкстры алгоритм Прима, который нам позволяет находить миностов.

Выберем начальную вершину $r$ -- корень нашего будущего миностова. Из этой вершины наш миностов будет <<проростать>>.
Далее все время будем жадно к нашему дереву присоединять вершины через минимальные рёбра и на выходе должен получиться наверное миностов (далее мы покажем, что это действительно так). Теперь к реализации.
Будем поддерживать наше дерево неявно, как множество пар <<вершина--родитель>>. 
Если мы построим алгоритм, что он будет нам возвращать такое множество, само дерево восстановить тривиально. Такое множество пар будем обозначать $T$.
Оно задается следующим образом:
\begin{equation*}
    T \overset{def}{=} \{ (\pi[v], v) \, | \, v \in V \}
\end{equation*}
Где $\pi$ -- отображение, которое вершине сопоставляет родителя в миностове.

Далее обстоит вопрос о том, как выбирать те самые минимальные рёбра. Будем поддерживать очередь вершин с приоритетами, где приоритет вершины $v$ -- это кратчайшее ребро до $T$ (то есть до некоторой вершины $T$).
Если такого ребра нет, то считаем, что это расстояние есть $+\infty$. Более формально положим:
\begin{equation*}
    k[v] \overset{def}{=} \min\limits_{(u, v) \in T} w(u, v)
\end{equation*}
Обратите внимание, что $k$ определяется относительного некоторой итерации алгоритма расширения остова.

Тогда результрующий алгоритм выглядит следующим образом (он очень похож и внешне и по механике на алгоритм Дейкстры):

\begin{algorithm}[H]
\caption{Алгоритм Прима}
\begin{algorithmic}[1]

\Function{PRIM}{$G\langle V, E \rangle, w$}
    \State $k \gets \{\}$
    \State $\pi \gets \{\}$

    \For{$v \in V$}
        \State $k[v] \gets +\infty$
        \State $\pi[v] \gets \varnothing$
    \EndFor

    \State $r \gets V[0]$
    \State $k[r] \gets 0$ \Comment{<<Расстояние>> от корня до миностова равно нулю, что логично взять базой}

    \State $Q \gets \texttt{MAKE\_QUEUE(}V, k\texttt{)}$ \Comment{Инициализируем приоритетную очередь на основе $k$}

    \While{$Q \neq \varnothing$}
        \State $v \gets Q\texttt{.EXTRACT\_MIN()}$

        \For{$(v, u) \in E$}
            \If{$u \in Q \land w(v, u) < k[u]$}
                \State $k[u] \gets w(v, u)$
                \State $\pi[u] \gets v$ \Comment{Обновляем родителя}
                \State $Q\texttt{.DECREASE\_KEY(}u, w(v, u)\texttt{)}$
            \EndIf
        \EndFor
    \EndWhile

    \State \Return $\{ (v, \pi[v]) \, | \, v \in V \}$
\EndFunction
\end{algorithmic}
\end{algorithm}

\St Функция \texttt{PRIM} возвращает миностов.
\begin{proof}

    Положим $\mathcal{V}_n$ -- множество всех достанных вершин на строчке 11 после $n$ итераций цикла. 
    Тогда положим $\mathcal{T}_n \overset{def}{=} \{ v \in \mathcal{V}_n \setminus \{ r \}:  (\pi[v], v) \}$ -- граф построенный на $n$-ом шаге.

    Покажем, что $\mathcal{T}_n$ на любом шаге является подграфом некоторого миностова по индукции.

    База $n = 0$. $\mathcal{T}_0 = \varnothing$ -- тривиально.
    
    Пусть мы сделали $n_0 - 1$ шагов цикла ($n_0 \geq 1$). Покажем корректность для $n_0$.

    Пусть мы достали из $Q$ вершину $v$ (далее считаем, что в $Q$ есть вершина $v$). 
    Рассмотрим разрез $(V \setminus Q, Q)$, если воспринимать $Q$ как множество вершин.
    Тогда заметим, что $(\pi[v], v)$ пересекает разрез.
    Более того, легко видеть, что данное ребро минимально по весу среди всех остальных рёбер пересекающих данный разрез.
    Далее, по предположению индукции, $\mathcal{T}_{n_0 - 1}$ -- подграф некоторого миностова.
    Значит данное ребро безопасно относительно $\mathcal{T}_{n_0 - 1}$, по лемме о безопасном ребре. 
    Значит, $\mathcal{T}_{n_0} = \mathcal{T}_{n_0 - 1} \cup \{(\pi[v], v)\}$ -- подграф некоторого миностова. Получили требуемое.

\end{proof}

Далее встаёт вопрос об асимптотике, которая напрямую зависит от реализации приоритетной очереди.

\St Если использовать бинарную пирамиду в качестве приоритетной очереди, то время работы есть $\O(|E| \log |V|)$
\begin{proof}
    
    Построение кучи занимает $\O(|V|)$. Далее \texttt{while} делает \texttt{EXTRACT\_MIN} $\O(|V|)$, значит они суммарно работают за $\O(|V| \log |V|)$.
    Далее, для каждого ребра (заметим, что каждое ребро просматривается ровно один раз) в худшем случае мы будем обновлять кучу за $\O(\log |V|)$.
    Значит, суммарно \texttt{DECREASE\_KEY} работает за $\O(|E| \log |V|)$.

    Итого, результирующее время выполнения $\O(|V| \log |V| + |E| \log |V|) = \O(|E| \log |V|)$.
    
\end{proof}

Казалось бы, получилось такая же временная сложность как и у алгоритма Крускала, но в отличие от него, алгоритм Прима можно ускорить.

\StBd Если использовать фибоначчиеву кучу, то время работы $\O(|E| + |V| \log |V|)$.

\begin{proof}

    Проводя аналогичные рассуждения, только учтя, что \texttt{EXTRACT\_MIN} работает за $\O^*(\log |V|)$, а \texttt{DECREASE\_KEY} за $\O^*(1)$.
    Получаем, что время выполнения $\O(|V| \log |V| + |E| \cdot 1) = \O(|V| \log |V| + |E|)$.

\end{proof}